%\chapter{Appendix}\label{chap: Appendix}

\chapter{Derivation of dispersive stress}\label{apx: derivation of dispersive stress}

This appendix derives the dispersive stress tensor introduced in \autoref{eq: dispersive stress tensor}.
The notation follows \autoref{par: dispersive stresses}.
An overline $\overline{ \phantom{u} }$ denotes a time average, angle brackets $\langle \phantom{u} \rangle$ denote a spatial average over the wall-parallel plane, a prime $ \phantom{u}'$ a fluctuation in time about the time average, and a tilde $\tilde{ \ }$ a fluctuation in space about the spatial average, i.e.\ $\tilde{u}_i = \overline{u_i} - \langle \overline{u_i}\rangle$.

\section{Analytical derivation}

The Reynolds stresses are defined as the time average of the product of the temporal velocity fluctuations
\begin{equation}\label{eq: reynolds stresses}
    R_{ij} = \overline{(u_i - \overline{u_i})(u_j - \overline{u_j})}
    = \overline{u_i' u_j'}
\end{equation}
The dispersive stress tensor follows from applying the same operation to the time-averaged velocity $\overline{u_i}$, but with the average taken in space rather than in time:
\begin{align}
    D_{ij}  &= \langle (\overline{u_i} - \langle \overline{u_i}\rangle)(\overline{u_j} - \langle \overline{u_j}\rangle) \rangle \notag\\
            &= \langle \overline{u_i}\,\overline{u_j} - \overline{u_i}\langle\overline{u_j}\rangle - \langle \overline{u_i}\rangle\overline{u_j} + \langle\overline{u_i}\rangle\langle\overline{u_j}\rangle \rangle \notag\\
    %\intertext{Since $\langle\overline{u_i}\rangle$ and $\langle\overline{u_j}\rangle$ no longer vary in space, they can be pulled out of the outer spatial average term by term:}
            &= \langle \overline{u_i}\,\overline{u_j} \rangle - \langle\overline{u_i}\rangle\langle\overline{u_j}\rangle - \langle\overline{u_i}\rangle\langle\overline{u_j}\rangle + \langle\overline{u_i}\rangle\langle\overline{u_j}\rangle \notag\\
            &= \langle \overline{u_i}\,\overline{u_j} \rangle - \langle\overline{u_i}\rangle\langle\overline{u_j}\rangle \notag\\
    %\intertext{Substituting $\overline{u_i} = \langle\overline{u_i}\rangle + \tilde{u}_i$ (and analogously for $\overline{u_j}$) and using that the spatial average of a spatial fluctuation vanishes, $\langle\tilde{u}_i\rangle = 0$, leaves only the cross term of the fluctuations:}
            &= \langle \tilde{u}_i\,\tilde{u}_j \rangle \label{eq: dispersive stress}
\end{align}
which recovers the definition of the dispersive stress tensor given in \autoref{eq: dispersive stress tensor}.

% \section{Evaluation in the post-processing script}

% The post-processing script computes the pointwise Reynolds stress $R_{ij} = \overline{u_i'u_j'}$ at every spatial location directly from the time-resolved velocity field, before any spatial averaging is applied.
% To recover the correctly double-averaged Reynolds stress $\langle R_{ij}\rangle$, the term $\langle \overline{u_i}\,\overline{u_j}\rangle - \langle\overline{u_i}\rangle\langle\overline{u_j}\rangle$ is added and subtracted, which makes the dispersive stress from \autoref{eq: dispersive stress} appear explicitly:
% \begin{align}
%     \langle R_{ij}\rangle   &= \langle R_{ij} + \overline{u_i}\,\overline{u_j} \rangle - \langle \overline{u_i} \rangle \langle\overline{u_j}\rangle \notag\\
%                             &= \langle \overline{u_i' u_j'} \rangle + \langle \overline{u_i}\,\overline{u_j} \rangle - \langle \overline{u_i} \rangle \langle\overline{u_j}\rangle \notag
%     \intertext{using \autoref{eq: dispersive stress} for the last two terms:}
%                             &= \langle \overline{u_i' u_j'} \rangle + \langle \tilde{u}_i\,\tilde{u}_j\rangle \notag\\
%                             &= \langle \overline{u_i' u_j'} \rangle + D_{ij} \label{eq: reynolds stress in PP}
% \end{align}
% i.e.\ the double-averaged Reynolds stress splits into the spatially averaged, pointwise Reynolds stress and the dispersive stress tensor.
% This is the identity used by the post-processing script: $\overline{u_i'u_j'}$ and $\overline{u_i}$ are computed directly from the simulation output, and $\langle\overline{u_i'u_j'}\rangle$ and $D_{ij}$ are then obtained from the spatial average and \autoref{eq: dispersive stress}, respectively.

% --- Original version, kept for reference ---
% The Reynolds stresses are defined as:
% \begin{equation}\label{eq: reynolds stresses}
%     R_{ij} = \overline{(u_i - \overline{u_i})(u_j - \overline{u_j})}
%     = \overline{u_i' u_j'}
% \end{equation}
%
% DErivation of Dispersive stresses.
% The * superscripts denoting spatial fluctiuations.
% \begin{equation}\label{eq: dispersive stress}
%     \begin{aligned}
%         D_{ij}  &= \langle (\overline{u_i} - \langle \overline{u_i}\rangle)(\overline{u_j} - \langle \overline{u_j}\rangle) \rangle\\
%                 &= \langle \overline{u_i}\,\overline{u_j} - \overline{u_i}\langle\overline{u_j}\rangle - \langle \overline{u_i}\rangle\overline{u_j} + \langle\overline{u_i}\rangle\langle\overline{u_j}\rangle \rangle\\
%                 &= \langle \overline{u_i}\,\overline{u_j} \rangle - \langle \overline{u_i}\langle\overline{u_j}\rangle \rangle - \langle \langle \overline{u_i}\rangle\overline{u_j} \rangle + \langle \langle\overline{u_i}\rangle\langle\overline{u_j}\rangle\rangle \\
%                 &= \langle \overline{u_i}\rangle \langle \overline{u_j}\rangle + \langle \overline{u_i}^* \, \overline{u_j}^*\rangle - \langle\overline{u_i}\rangle\langle\overline{u_j}\rangle - \langle\overline{u_i}\rangle\langle\overline{u_j}\rangle + \langle\overline{u_i}\rangle\langle\overline{u_j}\rangle \\
%                 &= \langle \overline{u_i}^* \, \overline{u_j}^*\rangle
%     \end{aligned}
% \end{equation}
%
% Derivation of the calculations of Reynolds and dispersive stresses, as done in the postprocessing script.
% \begin{equation}\label{eq: reynolds stress in PP}
%     \begin{aligned}
%         \langle R_{ij}\rangle   &= \langle R_{ij} + \overline{u_i}\,\overline{u_j} \rangle - \langle \overline{u_i} \rangle \langle\overline{u_j}\rangle  \\
%                                 &= \langle \overline{u_i' u_j'} + \overline{u_i}\,\overline{u_j} \rangle - \langle \overline{u_i} \rangle \langle\overline{u_j}\rangle \\
%                                 &= \langle \overline{u_i' u_j'} \rangle + \langle \overline{u_i}\,\overline{u_j} \rangle - \langle \overline{u_i} \rangle \langle\overline{u_j}\rangle \\
%                                 &= \langle \overline{u_i' u_j'} \rangle  + \langle \overline{u_i}\rangle \langle\overline{u_j}\rangle + \langle \overline{u_i}^* \,\overline{u_j}^*\rangle - \langle \overline{u_i} \rangle \langle\overline{u_j}\rangle \\
%                                 &= \langle \overline{u_i' u_j'} \rangle + \langle \overline{u_i}^* \,\overline{u_j}^*\rangle
%     \end{aligned}
% \end{equation}
%
% \begin{equation} \label{eq: dispersive stress in PP}
%     \begin{aligned}
%         D_{ij}  &= \langle\overline{u_i}\,\overline{u_j}\rangle - \langle\overline{u_i}\rangle\langle\overline{u_j}\rangle\\
%                 &= \langle \overline{u_i}\rangle \langle \overline{u_j}\rangle + \langle \overline{u_i}^* \, \overline{u_j}^*\rangle - \langle\overline{u_i}\rangle\langle\overline{u_j}\rangle\\
%                 &= \langle \overline{u_i}^* \, \overline{u_j}^*\rangle
%     \end{aligned}
% \end{equation}

\chapter{Energy spectra}\label{apx: Energy budgets}
\section{Energy spectra for $Re_b = 11700$}
\begin{figure}[h!]
    \centering
    \begin{subfigure}{0.495\textwidth}
        \subcaption{}
        \subimport{plots/Results/11700/}{11700_spectra_piro.pgf}
        \label{subfig: 11700 spectra piro}
    \end{subfigure}
    \hfill
        \begin{subfigure}{0.495\textwidth}
            \subcaption{}
            \subimport{plots/Results/11700/}{11700_spectra_kDown.pgf}
            \label{subfig: 11700 spectra kDown}
    \end{subfigure}
    \\[-20pt]
    \begin{subfigure}{0.495\textwidth}
        \subcaption{}
        \subimport{plots/Results/11700/}{11700_spectra_scaled.pgf}
        \label{subfig: 11700 spectra scaled}
    \end{subfigure}
    \centering
    \begin{subfigure}{\textwidth}
        \centering
        \hspace{0.9cm}
        \subimport{plots/Results/11700/}{3DSpectra_colorbar.pgf}
    \end{subfigure}
    \caption[Streamwise pre-multiplied energy spectra at $Re_b = 11700$]{Streamwise pre-multiplied Energy spectra at $Re_b = 11700$. Wall-normal position plotted over wavelength. Colour denoting energy content, with blue indicating low energy content, and yellow high energy content. Red marker indicates the position of maximal energy. Contour lines start at 0.25 with intervals of 0.5.~(a) $S1$.~(b) $S2$.~(c) $S3$}
    \label{fig: energy spectra 11700}
\end{figure}

\section{Energy spectra for $Re_b = 19000$}
\begin{figure}[h!]
    \centering
    \begin{subfigure}{0.495\textwidth}
        \subcaption{}
        \subimport{plots/Results/19000/}{19000_spectra_piro.pgf}
        \label{subfig: 19000 spectra piro}
    \end{subfigure}
    \hfill
        \begin{subfigure}{0.495\textwidth}
            \subcaption{}
            \subimport{plots/Results/19000/}{19000_spectra_kDown.pgf}
            \label{subfig: 19000 spectra kDown}
    \end{subfigure}
    \\[-20pt]
    \begin{subfigure}{0.495\textwidth}
        \subcaption{}
        \subimport{plots/Results/19000/}{19000_spectra_scaled.pgf}
        \label{subfig: 19000 spectra scaled}
    \end{subfigure}
    \centering
    \begin{subfigure}{\textwidth}
        \centering
        \hspace{0.9cm}
        \subimport{plots/Results/19000/}{3DSpectra_colorbar.pgf}
    \end{subfigure}
    \caption[Streamwise pre-multiplied energy spectra at $Re_b = 19000$]{Streamwise pre-multiplied Energy spectra at $Re_b = 19000$. Wall-normal position plotted over wavelength. Colour denoting energy content, with blue indicating low energy content, and yellow high energy content. Red marker indicates the position of maximal energy. Contour lines start at 0.25 with intervals of 0.5.~(a) $S1$.~(b) $S2$.~(c) $S3$}
    \label{fig: energy spectra 19000}
\end{figure}
\newpage
\section{Energy spectra for $S2$}

\begin{figure}[h!]
    \centering
    \begin{subfigure}{0.495\textwidth}
        \subcaption{}
        \subimport{plots/Results/comp_Reynolds/kDown/}{5300_spectra_kDown.pgf}
        \label{subfig: 5300 spectra kDown reyndolsComp}
    \end{subfigure}
    \hfill
        \begin{subfigure}{0.495\textwidth}
            \subcaption{}
            \subimport{plots/Results/comp_Reynolds/kDown/}{11700_spectra_kDown.pgf}
            \label{subfig: 11700 spectra kDown reynoldsComp}
    \end{subfigure}
    \\[-20pt]
    \begin{subfigure}{0.495\textwidth}
        \subcaption{}
        \subimport{plots/Results/comp_Reynolds/kDown/}{19000_spectra_kDown.pgf}
        \label{subfig: 19000 spectra kDwon reynoldsComp}
    \end{subfigure}
    \centering
    \begin{subfigure}{\textwidth}
        \centering
        \hspace{0.9cm}
        \subimport{plots/Results/comp_Reynolds/kDown/}{3DSpectra_colorbar.pgf}
    \end{subfigure}
    \caption[Streamwise pre-multiplied energy spectra for $S2$]{Streamwise pre-multiplied Energy spectra for $S2$ at the three Reynolds numbers. Wall-normal position plotted over wavelength. Colour denoting energy content, with blue indicating low energy content, and yellow high energy content. Red marker indicates the position of maximal energy. Contour lines start at 0.25 with intervals of 0.5.~(a) $Re_b = 5300$.~(b) $Re_b = 11700$.~(c) $Re_b = 19000$}
    \label{fig: energy spectra S2}
\end{figure}
\newpage
\section{Energy spectra for $S3$}
\begin{figure}[h!]
    \centering
    \begin{subfigure}{0.495\textwidth}
        \subcaption{}
        \subimport{plots/Results/comp_Reynolds/scaled/}{5300_spectra_scaled.pgf}
        \label{subfig: 5300 spectra scaled reyndolsComp}
    \end{subfigure}
    \hfill
        \begin{subfigure}{0.495\textwidth}
            \subcaption{}
            \subimport{plots/Results/comp_Reynolds/scaled/}{11700_spectra_scaled.pgf}
            \label{subfig: 11700 spectra scaled reynoldsComp}
    \end{subfigure}
    \\[-20pt]
    \begin{subfigure}{0.495\textwidth}
        \subcaption{}
        \subimport{plots/Results/comp_Reynolds/scaled/}{19000_spectra_scaled.pgf}
        \label{subfig: 19000 spectra scaled reynoldsComp}
    \end{subfigure}
    \centering
    \begin{subfigure}{\textwidth}
        \centering
        \hspace{0.9cm}
        \subimport{plots/Results/comp_Reynolds/scaled/}{3DSpectra_colorbar.pgf}
    \end{subfigure}
    \caption[Streamwise pre-multiplied energy spectra for $S3$]{Streamwise pre-multiplied Energy spectra for $S3$ at the three Reynolds numbers. Wall-normal position plotted over wavelength. Colour denoting energy content, with blue indicating low energy content, and yellow high energy content. Red marker indicates the position of maximal energy. Contour lines start at 0.25 with intervals of 0.5.~(a) $Re_b = 5300$.~(b) $Re_b = 11700$.~(c) $Re_b = 19000$}
    \label{fig: energy spectra S3}
\end{figure}



\chapter{Development of viscous drag with roughness height and streamwise slope}\label{apx: viscous drag to surface height}

\section{Surfaces $S2$ and $S3$ at $Re_b = 5300$}
\begin{figure}[h!]
    \centering
    \subimport{plots/Results/5300/}{5300Reb_kDowntauWall_over_height.pgf}
    \caption{Roughness height and viscous drag over streamwise position for  $S2$ at $5300 Re_b$}
\end{figure}
\begin{figure}[h!]
    \centering
    \hspace{0.9cm}
    \subimport{plots/Results/5300/}{5300Reb_kDown_tauWall_over_height_bySlope.pgf}
    \caption[Scatterplot of viscous drag over physical roughness height at $Re_b = 5300$ for $S2$]{Scatter plot of nondimensionalised viscous drag over roughness height. Colour denoting local streamwise slope. Data shown for $S2$ at $5300 Re_b$}
\end{figure}
\begin{figure}[h!]
    \centering
    \subimport{plots/Results/5300/}{5300Reb_scaledtauWall_over_height.pgf}
    \caption{Roughness height and viscous drag over streamwise position for  $S3$ at $5300 Re_b$}
\end{figure}
\begin{figure}[h!]
    \centering
    \hspace{0.9cm}
    \subimport{plots/Results/5300/}{5300Reb_scaled_tauWall_over_height_bySlope.pgf}
    \caption[Scatterplot of viscous drag over physical roughness height at $Re_b = 5300$ for $S3$]{Scatter plot of nondimensionalised viscous drag over roughness height. Colour denoting local streamwise slope. Data shown for $S3$ at $5300 Re_b$}
\end{figure}
\clearpage
\section{Surfaces $S2$ and $S3$ at $Re_b = 11700$}
\begin{figure}[h!]
    \centering
    \subimport{plots/Results/11700/}{11700Reb_kDowntauWall_over_height.pgf}
    \caption{Roughness height and viscous drag over streamwise position for  $S2$ at $11700 Re_b$}
\end{figure}
\begin{figure}[h!]
    \centering
    \hspace{0.9cm}
    \subimport{plots/Results/11700/}{11700Reb_kDown_tauWall_over_height_bySlope.pgf}
    \caption[Scatterplot of viscous drag over physical roughness height at $Re_b = 11700$ for $S2$]{Scatter plot of nondimensionalised viscous drag over roughness height. Colour denoting local streamwise slope. Data shown for $S2$ at $1700 Re_b$}
\end{figure}
\begin{figure}[h!]
    \centering
    \subimport{plots/Results/11700/}{11700Reb_scaledtauWall_over_height.pgf}
    \caption{Roughness height and viscous drag over streamwise position for  $S3$ at $11700 Re_b$}
\end{figure}
\begin{figure}[h!]
    \centering
    \hspace{0.9cm}
    \subimport{plots/Results/11700/}{11700Reb_scaled_tauWall_over_height_bySlope.pgf}
    \caption[Scatterplot of viscous drag over physical roughness height at $Re_b = 11700$ for $S3$]{Scatter plot of nondimensionalised viscous drag over roughness height. Colour denoting local streamwise slope. Data shown for $S3$ at $11700 Re_b$}
\end{figure}

\clearpage
\section{Surfaces $S2$ and $S3$ at $Re_b = 19000$}

\begin{figure}[h!]
    \centering
    \subimport{plots/Results/19000/}{19000Reb_kDowntauWall_over_height.pgf}
    \caption{Roughness height and viscous drag over streamwise position for  $S2$ at $19000 Re_b$}
\end{figure}
\begin{figure}[h!]
    \centering
    \hspace{0.9cm}
    \subimport{plots/Results/19000/}{19000Reb_kDown_tauWall_over_height_bySlope.pgf}
    \caption[Scatterplot of viscous drag over physical roughness height at $Re_b = 19000$ for $S2$]{Scatter plot of nondimensionalised viscous drag over roughness height. Colour denoting local streamwise slope. Data shown for $S2$ at $19000 Re_b$}
\end{figure}
\begin{figure}[h!]
    \centering
    \subimport{plots/Results/19000/}{19000Reb_scaledtauWall_over_height.pgf}
    \caption{Roughness height and viscous drag over streamwise position for  $S3$ at $19000 Re_b$}
\end{figure}
\begin{figure}[h!]
    \centering
    \hspace{0.9cm}
    \subimport{plots/Results/19000/}{19000Reb_scaled_tauWall_over_height_bySlope.pgf}
    \caption[Scatterplot of viscous drag over physical roughness height at $Re_b = 19000$ for $S3$]{Scatter plot of nondimensionalised viscous drag over roughness height. Colour denoting local streamwise slope. Data shown for $S3$ at $19000 Re_b$}
\end{figure}
